\chapter{Free Groups and Minimal Normal Subgroups}

\section{Free Groups}

\begin{definition}[Free group]
Let
\[
X_0=\{x_1,\dots,x_n\}
\]
be a set, and let
\[
X_0^{-1}=\{x_1^{-1},\dots,x_n^{-1}\}
\]
be a disjoint set of formal inverses. Set
\[
S=X_0\sqcup X_0^{-1}.
\]
A \textbf{word} is a finite sequence of elements of \(S\),
\[
w=a_1a_2\cdots a_r,\qquad a_i\in S.
\]
When \(r=0\), we obtain the empty word, denoted by \(1\).

A word is called a \textbf{reduced word} if it contains no adjacent pair
\[
x_i x_i^{-1}\quad \text{or}\quad x_i^{-1}x_i.
\]
Two words are called \textbf{equivalent} if one can be obtained from the other by finitely many insertions or deletions of
\[
x_i x_i^{-1},\qquad x_i^{-1}x_i.
\]

Write \([w]\) for the equivalence class of a word \(w\). Define multiplication by
\[
[u]\,[v]=[uv],
\]
where \(uv\) denotes the concatenation of the two words, followed by taking the equivalence class. All equivalence classes form a group, called the \textbf{free group} generated by
\[
x_1,\dots,x_n,
\]
denoted by
\[
F_n=\langle x_1,\dots,x_n\rangle.
\]
The number \(n\) is called the \textbf{rank} of the free group \(F_n\).
\end{definition}

\begin{remark}
The free group can also be defined equivalently as the set of all reduced words, with multiplication given by
\[
\text{concatenate, then reduce}.
\]
Hence every element of a free group can be written uniquely as a reduced word.
\end{remark}

\begin{example}[The free group of rank \(2\)]
Let
\[
F_2=\langle x_1,x_2\rangle.
\]
Then every nontrivial element \(g\in F_2\) can be written uniquely in a form where the generators alternate, for example
\[
g=x_1^{m_1}x_2^{m_2}x_1^{m_3}\cdots
\]
or
\[
g=x_2^{m_1}x_1^{m_2}x_2^{m_3}\cdots,
\]
where
\[
m_i\in \Z\setminus\{0\}.
\]
The key point is that adjacent generators cannot be formal inverses of one another; otherwise further reduction is possible.
\end{example}

\section{Generators and Relations}

A free group carries no relations other than those forced by the group axioms. Imposing additional relations on a free group yields many important groups.

\begin{example}[Dihedral groups]
We have the presentation
\[
\langle x_1,x_2\mid x_1^2=1,\ x_2^2=1\rangle.
\]
This is the \textbf{infinite dihedral group}.

Moreover,
\[
D_{2m}
=
\langle x_1,x_2\mid x_1^2=1,\ x_2^2=1,\ (x_1x_2)^m=1\rangle
\]
is the dihedral group of order \(2m\).
\end{example}

\begin{example}[Another presentation of the dihedral group]
Let
\[
a=x_1x_2,\qquad y=x_2.
\]
From
\[
x_1^2=x_2^2=1,\qquad (x_1x_2)^m=1
\]
we obtain
\[
a^m=1,\qquad y^2=1,\qquad yay^{-1}=a^{-1}.
\]
Therefore
\[
D_{2m}
=
\langle a,y\mid a^m=1,\ y^2=1,\ yay^{-1}=a^{-1}\rangle.
\]
\end{example}

\begin{definition}[Triangle group]
Let \(m,n,\ell\) be positive integers. The \textbf{triangle group} is defined by
\[
T_{m,n,\ell}
=
\langle x,y\mid x^m=1,\ y^n=1,\ (xy)^\ell=1\rangle.
\]
Geometrically, it is often viewed as the group generated by two elements \(x,y\) with
\[
|x|=m,\qquad |y|=n,\qquad |xy|=\ell.
\]
\end{definition}

Consider the angle sum
\[
\left(\frac1m+\frac1n+\frac1\ell\right)\pi.
\]
According to its size relative to \(\pi\), triangle groups fall into three types:
\[
\frac1m+\frac1n+\frac1\ell>1
\quad\Longleftrightarrow\quad
\text{spherical},
\]
\[
\frac1m+\frac1n+\frac1\ell=1
\quad\Longleftrightarrow\quad
\text{Euclidean},
\]
\[
\frac1m+\frac1n+\frac1\ell<1
\quad\Longleftrightarrow\quad
\text{hyperbolic}.
\]

Without loss of generality, assume
\[
m\leq n\leq \ell.
\]
If \(G\neq 1\), then \(\ell>1\).

\subsection*{Case 1: Spherical type}

If
\[
\frac1m+\frac1n+\frac1\ell>1,
\]
then the following cases may occur.

\begin{enumerate}
\item If \(m=1\), then \(x=1\) and the group reduces to the cyclic group generated by \(y\); when \(|y|=n\),
\[
G\cong C_n\cong \Z/n\Z.
\]

\item If \(m=2,\ n=2\), then
\[
G\cong D_{2\ell}.
\]

\item If \(m=2,\ n=3\), then we must have
\[
\ell=3,4,5.
\]
Further analysis yields:
\[
\begin{array}{c|c|c}
(m,n,\ell) & G & \text{Geometric meaning}\\ \hline
(2,3,3) & A_4 & \text{rotation group of the regular tetrahedron}\\
(2,3,4) & S_4 & \text{rotation group of the cube or regular octahedron}\\
(2,3,5) & A_5 & \text{rotation group of the regular dodecahedron or icosahedron}
\end{array}
\]
\end{enumerate}

\subsection*{Case 2: Euclidean type}

If
\[
\frac1m+\frac1n+\frac1\ell=1,
\]
then the possible triples are
\[
(m,n,\ell)=(3,3,3),\ (2,3,6),\ (2,4,4).
\]
These correspond to tilings of the Euclidean plane:
\[
\begin{array}{c|c}
(m,n,\ell) & \text{Corresponding tiling}\\ \hline
(3,3,3) & \text{equilateral triangle tiling}\\
(2,3,6) & \text{triangle/hexagon related tiling}\\
(2,4,4) & \text{square tiling}
\end{array}
\]
These groups are infinite and solvable.

\begin{remark}
In the Euclidean plane, only three regular polygons can tile the plane by themselves:
\[
\text{equilateral triangle},\qquad \text{square},\qquad \text{regular hexagon}.
\]
\end{remark}

\subsection*{Case 3: Hyperbolic type}

If
\[
\frac1m+\frac1n+\frac1\ell<1,
\]
then \(T_{m,n,\ell}\) is an infinite group and is generally nonsolvable. Such groups correspond to triangular tilings in hyperbolic geometry.

\section{An Example Involving Symmetric Groups}

\begin{example}
Let
\[
S_n=\langle (12),(12\cdots n)\rangle.
\]
Set
\[
x=(12),\qquad y=(12\cdots n).
\]
Then
\[
x^2=1,\qquad y^n=1.
\]
Moreover,
\[
xy=(23\cdots n),
\]
so
\[
(xy)^{n-1}=1.
\]
Thus we certainly have the relations
\[
x^2=y^n=(xy)^{n-1}=1.
\]

However, it is generally \emph{not} true that
\[
S_n
=
\langle x,y\mid x^2=1,\ y^n=1,\ (xy)^{n-1}=1\rangle;
\]
more relations are needed. Otherwise, when \(n\) is sufficiently large,
\[
\frac12+\frac1n+\frac1{n-1}<1,
\]
and by the triangle-group discussion, the group given by this presentation would be infinite, contradicting the finiteness of \(S_n\).
\end{example}

\section{Finitely Generated Groups as Quotients of Free Groups}

\begin{theorem}
Every finitely generated group is a homomorphic image of some free group.
\end{theorem}

\begin{proof}
Let
\[
G=\langle g_1,\dots,g_n\rangle.
\]
Let
\[
F_n=\langle x_1,\dots,x_n\rangle
\]
be the free group of rank \(n\). Define a map
\[
\varphi\colon F_n\longrightarrow G
\]
by
\[
\varphi(x_i)=g_i,\qquad
\varphi(x_i^{-1})=g_i^{-1}.
\]
Since elements of the free group are reduced words, \(\varphi\) extends naturally to a group homomorphism.

Because \(g_1,\dots,g_n\) generate \(G\), the map \(\varphi\) is surjective. Let
\[
N=\ker\varphi.
\]
By the first isomorphism theorem,
\[
F_n/N\cong G.
\]
Hence \(G\) is a homomorphic image of the free group \(F_n\).
\end{proof}

\begin{definition}[Defining relations]
Let
\[
G\cong F_n/N.
\]
Then \(x_1,\dots,x_n\) are called generators of \(G\), and the elements of \(N\) are called relations.

If \(N\), as a normal subgroup, is generated by a set \(R\subseteq F_n\), that is,
\[
N=\langle R\rangle,
\]
we write
\[
G=\langle x_1,\dots,x_n\mid r=1,\ r\in R\rangle.
\]
\end{definition}

\section{Minimal Normal Subgroups}

Let
\[
\Omega=\{1,2,\dots,n\},\qquad \operatorname{Sym}(\Omega)=S_n.
\]

\begin{definition}[Minimal normal subgroup]
Let \(G\) be a group. A nontrivial normal subgroup \(M\lhd G\) is called a \textbf{minimal normal subgroup} of \(G\) if there is no nontrivial normal subgroup \(N\lhd G\) with
\[
1<N<M.
\]
\end{definition}

\begin{definition}[Transitive and regular actions]
Let \(G\leq \operatorname{Sym}(\Omega)\).

If for all \(\alpha,\beta\in\Omega\) there exists \(g\in G\) such that
\[
\alpha^g=\beta,
\]
then \(G\) is said to act \textbf{transitively} on \(\Omega\).

If \(G\) is transitive and the point stabilizer of every point is trivial, that is,
\[
G_\alpha=1,
\]
then \(G\) is said to act \textbf{regularly} on \(\Omega\).
\end{definition}

\begin{remark}
If \(H\leq S_n\) is not transitive on \(\Omega\), then it preserves a nontrivial orbit decomposition. In particular, if there are two orbits of sizes \(k\) and \(n-k\), then
\[
H\leq S_k\times S_{n-k}.
\]
\end{remark}

\begin{definition}[Quasi-primitive permutation group]
Let \(G\leq \operatorname{Sym}(\Omega)\). If every nontrivial normal subgroup of \(G\) acts transitively on \(\Omega\), then \(G\) is called a \textbf{quasi-primitive permutation group}, or a \textbf{quasi-primitive group}.
\end{definition}

\begin{example}[A counterexample]
Consider the dihedral group in its natural action,
\[
G=D_{2n}=\langle a,b\rangle\leq S_n,
\]
where
\[
a=(12\cdots n)
\]
is a rotation and
\[
b=(1\,n)(2\,n-1)\cdots
\]
is a reflection.

If \(n\) is even, then
\[
N=\langle a^2\rangle=\{1,a^2,\dots,a^{n-2}\}
\]
is a normal subgroup of \(G\), and
\[
|N|=\frac n2.
\]
It has two orbits on \(\Omega\), namely the odd and even points, so it is not transitive. Hence in this case \(D_{2n}\) is not quasi-primitive.
\end{example}

\begin{proposition}
Let \(G\leq \operatorname{Sym}(\Omega)\) be a quasi-primitive permutation group, and let \(M\) be a minimal normal subgroup of \(G\). Then:

\begin{enumerate}
\item \(M\) acts transitively on \(\Omega\);
\item \(M\) is a direct product of isomorphic simple groups, that is,
\[
M\cong T_1\times \cdots \times T_r,
\]
where
\[
r\geq 1,\qquad T_1\cong \cdots \cong T_r\cong T,
\]
and \(T\) is a simple group.
\end{enumerate}
\end{proposition}

\begin{proof}
Since \(G\) is quasi-primitive and \(M\lhd G\) is nontrivial, \(M\) must be transitive by definition.

Since \(M\) is a minimal normal subgroup of \(G\), it has no nontrivial proper characteristic subgroups. Otherwise, if \(K\) is a nontrivial proper characteristic subgroup of \(M\), then because \(M\lhd G\), we have
\[
K\lhd G,
\]
contradicting the minimality of \(M\).

Therefore \(M\) is a finite characteristically simple group. Every finite characteristically simple group is a direct product of isomorphic simple groups, so
\[
M\cong T_1\times \cdots \times T_r,
\]
where the \(T_i\) are pairwise isomorphic and simple.
\end{proof}

\begin{remark}
If \(G\) is an arbitrary permutation group, a minimal normal subgroup need not be transitive. For example,
\[
G=\{1,(12)(34),(13)(24),(14)(23)\}\leq S_4.
\]
This is the Klein four-group. Since \(G\) is abelian, every subgroup is normal. Let
\[
A=\{1,(12)(34)\}\lhd G.
\]
Then \(A\) is a minimal normal subgroup of \(G\), but it is not transitive on \(\{1,2,3,4\}\).
\end{remark}

\begin{lemma}
Let \(M\lhd G\) and suppose \(M\) acts transitively on \(\Omega\). Then for every \(\omega\in\Omega\),
\[
G=MG_\omega,
\]
where
\[
G_\omega=\{g\in G\mid \omega^g=\omega\}
\]
is the stabilizer of \(\omega\).
\end{lemma}

\begin{proof}
Take any \(g\in G\) and set
\[
\omega'=\omega^g.
\]
Since \(M\) is transitive, there exists \(m\in M\) such that
\[
\omega^m=\omega'.
\]
Then
\[
\omega^{gm^{-1}}
=
(\omega^g)^{m^{-1}}
=
(\omega^m)^{m^{-1}}
=
\omega.
\]
Therefore
\[
gm^{-1}\in G_\omega.
\]
Hence
\[
g\in G_\omega M.
\]
Since \(M\lhd G\), we have
\[
G_\omega M=MG_\omega.
\]
Thus
\[
G=MG_\omega.
\]
\end{proof}

\section{Minimal Normal Subgroups of Quasi-Primitive Groups}

\begin{lemma}
Let \(G\leq \operatorname{Sym}(\Omega)\) be a quasi-primitive permutation group. Then \(G\) has at most two minimal normal subgroups.

If \(M,N\) are two distinct minimal normal subgroups of \(G\), then
\[
M\cong N,
\]
and both \(M\) and \(N\) are nonabelian regular groups.
\end{lemma}

\begin{proof}
Let \(M,N\) be two distinct minimal normal subgroups. Since
\[
M\cap N\lhd G,
\]
and since \(M,N\) are minimal normal and distinct, we have
\[
M\cap N=1.
\]
Moreover,
\[
[M,N]\leq M\cap N,
\]
so
\[
[M,N]=1.
\]
Hence \(M\) and \(N\) commute pointwise, and
\[
MN=M\times N\lhd G.
\]

Since \(G\) is quasi-primitive, both \(M\) and \(N\) are nontrivial normal subgroups, so both act transitively on \(\Omega\).

We now show that \(M\) is regular. Take any \(\omega,\omega'\in\Omega\). Since \(N\) is transitive, there exists \(x\in N\) such that
\[
\omega^x=\omega'.
\]
Then
\[
M_{\omega'}
=
M_{\omega^x}
=
x^{-1}M_\omega x.
\]
But \(M\) and \(N\) commute pointwise, so
\[
x^{-1}M_\omega x=M_\omega.
\]
Hence
\[
M_{\omega'}=M_\omega
\]
for every \(\omega'\in\Omega\).

Therefore, if \(m\in M_\omega\), then \(m\) fixes every point \(\omega'\in\Omega\). Since the action of \(G\leq \operatorname{Sym}(\Omega)\) is faithful, we get
\[
m=1.
\]
Thus
\[
M_\omega=1.
\]
So \(M\) is regular. By the same argument, \(N\) is regular.

Next we show that \(M\) and \(N\) are nonabelian. Suppose for contradiction that \(M\) is abelian. Choose \(1\neq x\in N\), and let
\[
H=\langle M,x\rangle.
\]
Since \(M\) and \(N\) commute pointwise, \(H\) is abelian. Since \(M\) is transitive, \(H\) is also transitive.

An abelian transitive permutation group must be regular: if \(A\) is abelian and transitive, then
\[
A_{\omega^a}=a^{-1}A_\omega a=A_\omega
\]
for every \(a\in A\), so a point stabilizer fixes all points; by faithfulness, the stabilizer is trivial.

Hence \(H\) is regular. Thus
\[
|H|=|\Omega|.
\]
But \(M\) is also regular, so
\[
|M|=|\Omega|.
\]
Since \(x\notin M\), we have
\[
M<H,
\]
and therefore
\[
|H|>|M|,
\]
a contradiction. So \(M\) is nonabelian. By the same argument, \(N\) is nonabelian.

We now prove that \(M\cong N\). Let
\[
X=MN=M\times N.
\]
Since \(M\) is transitive, \(X\) is transitive. Fix a point \(\omega\in\Omega\) and consider the stabilizer \(X_\omega\).

Since \(M\) and \(N\) are regular,
\[
|M|=|N|=|\Omega|.
\]
Thus
\[
|X_\omega|
=
\frac{|X|}{|\Omega|}
=
\frac{|M||N|}{|\Omega|}
=
|\Omega|.
\]
Consider the natural projection
\[
\pi_M\colon X_\omega\longrightarrow M.
\]
Its kernel is
\[
\ker\pi_M=X_\omega\cap N=N_\omega=1,
\]
so \(\pi_M\) is injective. Also,
\[
|X_\omega|=|M|,
\]
so \(\pi_M\) is bijective.

Therefore, for each \(m\in M\), there is a unique \(\psi(m)\in N\) such that
\[
m\psi(m)\in X_\omega.
\]
This defines a map
\[
\psi\colon M\longrightarrow N.
\]
Since \(X_\omega\) is a subgroup, if
\[
m_1\psi(m_1),\quad m_2\psi(m_2)\in X_\omega,
\]
then
\[
m_1m_2\psi(m_1)\psi(m_2)\in X_\omega.
\]
By uniqueness,
\[
\psi(m_1m_2)=\psi(m_1)\psi(m_2).
\]
So \(\psi\) is a homomorphism. Similarly, the projection onto \(N\) is bijective, hence \(\psi\) is an isomorphism. Therefore
\[
M\cong N.
\]

Finally, we show that \(G\) has at most two minimal normal subgroups. Suppose \(L\) is a third minimal normal subgroup with
\[
L\neq M.
\]
By the preceding discussion, \(L\) is also regular and commutes pointwise with \(M\), so
\[
L\leq C_{\operatorname{Sym}(\Omega)}(M).
\]

On the other hand, if \(M\) is a regulargroup, then each element of \(C_{\operatorname{Sym}(\Omega)}(M)\) is determined uniquely by its value at one point, so
\[
|C_{\operatorname{Sym}(\Omega)}(M)|\leq |\Omega|.
\]
Since \(N\) is regular and commutes pointwise with \(M\),
\[
N\leq C_{\operatorname{Sym}(\Omega)}(M),
\qquad
|N|=|\Omega|.
\]
Therefore
\[
C_{\operatorname{Sym}(\Omega)}(M)=N.
\]
Similarly,
\[
L\leq C_{\operatorname{Sym}(\Omega)}(M)=N.
\]
Since \(L\) and \(N\) are both nontrivial minimal normal subgroups,
\[
L=N.
\]
Hence \(G\) has at most two minimal normal subgroups.
\end{proof}

\begin{corollary}
    Let \(G\leq \operatorname{Sym}(\Omega)\) be transitive. Then
    \[
        C_{\operatorname{Sym}(\Omega)}(G)
    \]
    is semiregular on \(\Omega\).
    \end{corollary}
    
    \begin{proof}
    Let
    \[
        C=C_{\operatorname{Sym}(\Omega)}(G).
    \]
    We prove that
    \[
        C_\omega=1
        \qquad
        \text{for every } \omega\in\Omega.
    \]
    
    Fix \(\omega\in\Omega\), and let \(c\in C_\omega\). Thus
    \[
        \omega^c=\omega.
    \]
    Let \(\omega'\in\Omega\). Since \(G\) is transitive on \(\Omega\), there exists
    \(g\in G\) such that
    \[
        \omega^g=\omega'.
    \]
    Since \(c\in C_{\operatorname{Sym}(\Omega)}(G)\), we have
    \[
        cg=gc.
    \]
    Therefore
    \[
        (\omega')^c
        =
        (\omega^g)^c
        =
        \omega^{gc}
        =
        \omega^{cg}
        =
        (\omega^c)^g
        =
        \omega^g
        =
        \omega'.
    \]
    Thus \(c\) fixes every point of \(\Omega\). Hence \(c=1\). Therefore
    \[
        C_\omega=1.
    \]
    Since \(\omega\in\Omega\) was arbitrary, \(C\) is semiregular on \(\Omega\).
    \end{proof}

    \begin{definition}[Semiregular permutation group]
        Let \(L\leq \operatorname{Sym}(\Omega)\). We say that \(L\) is
        \textbf{semiregular} on \(\Omega\) if
        \[
            L_\omega=1
            \qquad
            \text{for every } \omega\in\Omega,
        \]
        where
        \[
            L_\omega=\{\ell\in L\mid \omega^\ell=\omega\}
        \]
        is the stabilizer of \(\omega\) in \(L\).
        
        Equivalently, \(L\) is semiregular if no nonidentity element of \(L\) fixes any point of
        \(\Omega\); that is,
        \[
            1\neq \ell\in L
            \quad\Longrightarrow\quad
            \omega^\ell\neq \omega
            \qquad
            \text{for all } \omega\in\Omega.
        \]
        \end{definition}


